In this section we introduce an extension to the semantics of CEL, namely we introduce the \textit{joins} operator to combine complex events that agree on their shared variables. 
We then extend the formal computational model for evaluating CEL formulas and prove its correctness. 

Given two complex events $C_1 = (i_1, j_1, \mu_1)$ and $C_2 = (i_2, j_2, \mu_2)$, we say that $C_1$ and $C_2$ are \emph{compatible}, denoted by $C_1 \sim C_2$, if $C_1(X) = C_2(X)$ for every $X \in \text{dom}(C_1) \cap \text{dom}(C_2)$. We present then the \emph{joins operator} for CEL with the following definition, where $\text{dom}(\varphi_1 \joins \varphi_2) = \text{dom}(\varphi_1) \cup \text{dom}(\varphi_2)$:
\[ 
\begin{aligned}
	\sem{\varphi_1 \joins \varphi_2}(\Stream)
	&= \{C_1 \cup C_2 \mid C_1 \in \sem{\varphi_1}(\Stream) 
	\land C_2 \in \sem{\varphi_2}(\Stream) 
	\land C_1 \sim C_2\}
\end{aligned}
\]

\begin{example} Suppose we have sensors in a park monitoring humidity ($H$), temperature ($T$), and smoke ($S$). Consider a stream $\Stream = e_1 e_2 e_3$ where the event types are $\type(e_1) = H$, $\type(e_2) = T$, and $\type(e_3) = S$. 
	
	Let $\varphi_1 = H ; T$ be a formula detecting a humidity reading followed by a temperature reading, with $\text{dom}(\varphi_1) = \{H, T\}$, and $\varphi_2 = T ; S$ be a formula detecting a temperature reading followed by a smoke reading, with $\text{dom}(\varphi_2) = \{T, S\}$.
	\begin{itemize}
		\item The set $\sem{\varphi_1}(\Stream)$ contains a complex event $C_1$ with $\cinterval(C_1) = [1..2]$, $C_1(H) = \{1\}$, and $C_1(T) = \{2\}$.
		\item The set $\sem{\varphi_2}(\Stream)$ contains a complex event $C_2$ with $\cinterval(C_2) = [2..3]$, $C_2(T) = \{2\}$, and $C_2(S) = \{3\}$.
	\end{itemize}
	
	The shared domain between both events is $\text{dom}(C_1) \cap \text{dom}(C_2) = \{T\}$. Since $C_1(T) = C_2(T) = \{2\}$, the compatibility condition holds, and the complex event $C = C_1 \cup C_2$ is in the result of $\sem{\varphi_1 \joins \varphi_2}(\Stream)$ with $\cinterval(C) = [1..3]$, $C(H) = \{1\}$, $C(T) = \{2\}$, and $C(S) = \{3\}$. This resulting complex event uses the shared temperature reading at position $2$ as a bridge connecting the initial humidity event to the subsequent smoke event.
\end{example}

We denote by $\celplus{\joinssimb}$ the \emph{extension of CEL} with the joins operation. We can prove the following result:

\begin{theorem} \label{theo:joins}
	For every $\celplus{\joinssimb}$ formula $\varphi$ there exists a CEA $\cA_\varphi$ such that $\sem{\varphi}(\Stream) = \sem{\cA_\varphi}(\Stream)$ for every stream $\Stream$.
\end{theorem}

\begin{proof}
	We already know that for every (standard) CEL operator we can construct an equivalent CEA. So, we concentrate next in the construction of a CEA for the joins operator. 
	Let $\varphi_1$ and $\varphi_2$ be formulas in CEL extended with the joins operator. We show by induction that there exists a CEA $\cAjoins$ such that:
	\[
	\sem{\varphi_1 \joins \varphi_2}(\Stream) = \sem{\cAjoins}(\Stream)
	\]
	Assume then that there exist CEAs that satisfy the previous property for $\varphi_1$ and $\varphi_2$, called $\cA_{\varphi_1} = (Q_1, \SETP_1, \SETX_1, \Delta_1, q_0, F_1)$ and $\cA_{\varphi_2} = (Q_2, \SETP_2, \SETX_2, \Delta_2, p_0, F_2)$, respectively. 
	Then the construction for the joins operator is as follows. 
	
	\paragraph{$\sem{\varphi_1 \joins \varphi_2}(\Stream) \subseteq \sem{\cAjoins}(\Stream)$}
	\vspace{4cm}
	
	\paragraph{$\sem{\cAjoins}(\Stream) \subseteq \sem{\varphi_1 \joins \varphi_2}(\Stream)$}
	\vspace{4cm}
	
\end{proof}